\worksheettitle
  {M05\(\star\). Проверяем полноту}
  {\modulemark{M05\(\star\)} Усложнение без полного перестроения списка}

\studentnameblock

\begin{warningbox}[Какой список полный?]
  Из цифр 1, 3, 8 составляли трёхзначные числа без повторений.

  Список А: $138,183,318,381,813,831$.

  Список Б: $138,183,318,381,813$.

  Не переписывая варианты заново, объясните, какой список полный:
  \answerblank[120mm].
\end{warningbox}

\begin{practicebox}[Найдите пропуск по структуре]
  Цифры 0, 2, 5, числа трёхзначные, без повторений:
  $205,250,502$.

  Какая группа по первой цифре неполна? \answerblank[25mm].
  Какого числа не хватает? \answerblank[25mm].
  Объяснение без полного перебора: \answerblank[105mm].
\end{practicebox}

\begin{practicebox}[Создайте проверяемый список]
  Выберите три разные ненулевые цифры. Составьте все трёхзначные числа без
  повторений и расположите их группами по первой цифре.

  Цифры: \answerblank[45mm].

  Группа 1: \answerblank[80mm].

  Группа 2: \answerblank[80mm].

  Группа 3: \answerblank[80mm].

  Одно число, которое можно удалить как ловушку: \answerblank[35mm].
  Как быстро обнаружить пропуск? \answerblank[90mm].
\end{practicebox}
